\section{Introduction}


The advent of large-scale quantum computers poses an existential threat to current blockchain systems. Shor's algorithm can break RSA and ECDSA in polynomial time, compromising the security of $>$99\% of deployed blockchains \cite{shor1997}. While quantum computers with sufficient qubits remain years away, blockchain security must be \emph{proactive}, not reactive.

\subsection{The Quantum Threat}

Classical blockchain security relies on computational hardness assumptions:
\begin{itemize}
    \item \textbf{ECDSA signatures}: Secure against classical computers (solving discrete log requires $O(2^{128})$ operations)
    \item \textbf{Quantum vulnerability}: Shor's algorithm solves discrete log in $O(n^3)$ time
    \item \textbf{Timeline}: NIST estimates quantum threat by 2030-2035 \cite{nist2022}
\end{itemize}

\textbf{Harvest-now-decrypt-later attacks}: Adversaries can store encrypted blockchain data today and decrypt it once quantum computers become available, retroactively compromising all transactions.

\subsection{Our Contribution}

Lux Quantum Consensus introduces:

\begin{enumerate}
    \item \textbf{Post-quantum signature schemes}: Dilithium (3,293-byte signatures) for validators, SPHINCS+ for long-term security
    \item \textbf{Lattice-based threshold signatures}: Distributed key generation immune to quantum attacks
    \item \textbf{Quantum-safe finality}: Hybrid consensus combining Wave family with quantum-resistant cryptography
    \item \textbf{Zero-knowledge post-quantum proofs}: zk-STARKs for privacy-preserving validation
    \item \textbf{Backward compatibility}: Gradual migration from ECDSA via hybrid signatures
\end{enumerate}

\textbf{Performance}: We maintain Lux's core properties:
\begin{itemize}
    \item Sub-2-second finality
    \item 50,000+ TPS throughput
    \item Byzantine fault tolerance (33\% adversarial threshold)
    \item Cross-chain interoperability
\end{itemize}

